% ============================================================
% Section 134: 电离杂质散射与电导率温度依赖
% ============================================================
\section{半导体物理：电离杂质散射与电导率的温度依赖}
\label{sec:ionized-impurity-scattering}

\begin{modelbox}[题目：电离杂质散射主导下的电导率温度关系]
求半导体中电导率 $\sigma=ne^2\tau/m$ 对温度的依赖关系（其中 $\tau$ 为弛豫时间）。该半导体处在使其中自由载流子浓度为常数的温度下，并且散射机制是由于少量而恒定的杂质电荷所引起的卢瑟福散射。
\end{modelbox}

\begin{tipbox}[考点快速分析]
\begin{itemize}
    \item \textbf{卢瑟福散射}：带电载流子与电离杂质中心的库仑散射
    \item \textbf{弛豫时间}：$1/\tau\propto N_Iv^{-3}$，康威尔-韦斯可夫/布鲁克斯-赫林公式
    \item \textbf{能量均分}：$\frac{1}{2}m^*v_{\text{th}}^2=\frac{3}{2}k_BT\implies v\propto T^{1/2}$
    \item \textbf{电导率}：$\sigma=ne^2\tau/m^*\propto\tau\propto T^{3/2}$（当 $n$、$N_I$ 恒定）
\end{itemize}
\end{tipbox}

\begin{derivationbox}[标准解题过程]
\textbf{1. 卢瑟福散射机制与弛豫时间}

电离杂质散射的微分散射截面为
\[
\sigma(\theta)=\left(\frac{Ze^2}{8\pi\varepsilon_0\varepsilon_rm^*v^2}\right)^2\csc^4\frac{\theta}{2}.
\]
散射概率（弛豫时间的倒数）由考虑动量损失的积分给出：
\[
\frac{1}{\tau}=N_Iv\int(1-\cos\theta)\sigma(\theta)\,d\Omega,
\]
其中 $N_I$ 为电离杂质浓度。代入并积分后得到
\[
\frac{1}{\tau}=\frac{N_I Z^2e^4}{8\pi(\varepsilon_0\varepsilon_r)^2m^{*2}v^3}\ln\left[1+\left(\frac{2\varepsilon_0\varepsilon_rm^*v^2}{Ze^2N_I^{1/3}}\right)^2\right].
\]
对数项随 $v$ 和 $N_I$ 变化缓慢，在近似分析中可视为常数。因此
\[
\frac{1}{\tau}\propto N_Iv^{-3}.
\]

\textbf{2. 速度与温度的关系}

在非简并半导体中，由能量均分定理
\[
\frac{1}{2}m^*v_{\text{th}}^2=\frac{3}{2}k_BT\implies v_{\text{th}}\propto T^{1/2}.
\]
因此
\[
\frac{1}{\tau}\propto N_IT^{-3/2}.
\]

\textbf{3. 电导率的温度依赖关系}

电导率公式 $\sigma=ne^2\tau/m^*$。已知载流子浓度 $n$ 为常数（饱和电离区），$N_I$ 恒定，故
\begin{equation}
\boxed{\sigma\propto\tau\propto T^{3/2}.}
\end{equation}

\textbf{物理意义}：温度升高，载流子热运动速度增大，快速通过杂质中心附近，受库仑散射的时间缩短，散射概率降低，迁移率增大，从而电导率上升。

\textbf{与晶格散射的对比}：声子散射主导时 $\mu\propto T^{-3/2}$，电导率随温度升高而下降。因此半导体迁移率-温度曲线常呈现：低温区电离杂质散射主导（$\mu\propto T^{3/2}$），高温区声子散射主导（$\mu\propto T^{-3/2}$），中间存在迁移率峰值。

\textbf{适用范围}：本题假设载流子浓度恒定，对应于非本征饱和区。当温度进一步升高进入本征区时，$n\propto T^{3/2}e^{-E_g/2k_BT}$，本征激发的指数项将主导电导率，使其急剧上升。
\end{derivationbox}

\begin{formulabox}[答案汇总]
\begin{center}
\begin{tabular}{ll}
\toprule
\textbf{结论} & \textbf{结果} \\
\midrule
弛豫时间与温度 & $\tau\propto T^{3/2}$ \\
电导率与温度 & $\sigma\propto T^{3/2}$ \\
物理原因 & 温度升高 $\to$ 速度增大 $\to$ 散射减弱 $\to$ 迁移率增大 \\
\bottomrule
\end{tabular}
\end{center}
\end{formulabox}

\begin{insightbox}[物理图像]
\begin{itemize}
    \item 电离杂质散射的弛豫时间 $\tau\propto v^3\propto T^{3/2}$。恒定载流子浓度下，电导率的温度依赖完全由迁移率决定，即 $\sigma\propto\mu\propto T^{3/2}$。
    \item 这一特征是区分半导体中散射机制（杂质散射 vs 晶格散射）的重要实验判据。通过测量迁移率-温度关系，可以判断样品中 dominant 的散射机制。
    \item 实际半导体中，低温区通常以电离杂质散射为主（高纯样品例外），高温区以声子散射为主，中间存在一个迁移率最大值。
\end{itemize}
\end{insightbox}
